\section{The moving lower endpoint \texorpdfstring{$T=1/L$}{T=1/L}}
\label{sec:moving-endpoint}
This section integrates the closed moving-endpoint argument of
\cite{rec20,rec20review}. It uses the earlier fixed-domain machinery
with an explicit new endpoint ledger; the norms and proofs on
\([1,2]\) retain their original meaning. Put
\begin{equation}
 \begin{gathered}
 R=r+1,\quad Y=b+1,\quad W=R+Y,\quad
 \nu=\lceil W\rceil+1=\lceil r+b+2\rceil+1,\\
 L=\log(X/(4\pi)),\quad p=\lfloor L/2\rfloor,\quad
 K=2p,\quad V=8\sqrt2L.
 \end{gathered}
 \label{eq:moving-data}
\end{equation}
Every descriptor and the lower endpoint are frozen under \(z,T\)
differentiation. The original unit buffers remain, as do both arms,
every \(n\le N\), every \(C_0,\ldots,C_K\), the exact adjacent-cutoff
transfer, and the bounded real completion integral in \eqref{eq:model}.

\begin{definition}[Moving-domain channel norm]
\label{def:moving-channel-norm}
For independently supplied channels
\(\mathbf E=(E_0,E_1,E_2,E_3,E_T)\), define
\begin{equation}
 \lVert\mathbf E\rVert_{\mathrm o;X;\mathrm{mov}}^{(3,1)}
 =\sup_{\substack{L^{-1}\le T\le2\\z\in\Omega_X(r,b)}}
 \frac{\max_{0\le j\le3}|E_j(T,z)|+|E_T(T,z)|}
      {\calE_{\mathrm o}(-T,z)}.
 \label{eq:moving-norm}
\end{equation}
For an actual smooth field \(E\), use its derivative tuple in this
notation. Replacing the denominator by \(\abs{G(z)}\) gives the
moving natural norm, and by \(g=G(X)\) the moving center norm.
The norm \(\lVert\cdot\rVert_{\mathrm o;X;1,2}^{(3,1)}\) is unchanged.
\end{definition}

\begin{theorem}[Moving-endpoint canonical delivery]
\label{thm:moving-endpoint-delivery}
Fix positive dyadic chart radii \(r,b\). With the canonical normalization
\eqref{eq:canonical}, \eqref{eq:moving-data} and the unchanged calibration
\eqref{eq:explicit-calibration}, a sufficient all-height onset is
\begin{equation}
 L\ge L_{\mathrm{mov}}=1024\nu^2.
 \qquad X\ge X_{\mathrm{mov}}=2^{2048\nu^2+8}
 \quad\hbox{is a stronger sufficient condition}.
 \label{eq:moving-onset}
\end{equation}
The exact completed model has moving outer comparison error
at most \(\tfrac12\e^{-L^2/4}\). On the strictly witnessed rational
queries, or the literal endpoint tag, specified in
\cref{subsec:moving-input}, finite constructions deliver independent
field, spatial-through-three and first-time channels with
\begin{equation}
 \bigl\lVert
   (\widehat h_j-\partial_z^jh)_{j=0}^3,
                  \widehat h_T-\partial_Th
 \bigr\rVert_{\mathrm o;X;\mathrm{mov}}^{(3,1)}
 \le\e^{-L^2/4}.
 \label{eq:moving-delivery}
\end{equation}
Here the supremum is over accepted queries in the displayed domain;
the exact mathematical comparison includes its closed lower endpoint.
Sufficient outer quadrature order per branch is \(O_{r,b}(L^3)\), while
sufficient delivered precision is
\begin{equation}
 P=\frac{3L^2}{8\log2}+O_{r,b}(L\log(L+2)).
 \label{eq:moving-precision}
\end{equation}
Provider-internal work and witness search costs are separate, as in
\cref{subsec:moving-cost}. On \([1,2]\), the sharper sufficient onset
\(L\ge64\nu\) in \cref{thm:linear-onset,thm:effective} remains valid.
\end{theorem}
The theorem is proved by the following kernel, truncation, delivery and
input arguments. It concerns an analytical finite construction; complete
canonical evaluator software and its certification remain open.

\subsection{Kernel normalization and integrated moments}
\label{subsec:moving-moments}
For \(z=x+\ii y\), \(v=u-x\), the exact modulus is
\begin{equation}
 |k_T(z-u)|=(4\pi T)^{-1/2}
       \exp(y^2/(4T))\exp(-v^2/(4T)).
 \label{eq:moving-kernel}
\end{equation}
Thus \(T\ge1/L\) costs normalization at most
\(\sqrt L/\sqrt{4\pi}\), transverse growth at most
\(\e^{Y^2L/4}\), and actual inverse-time derivative factors.
For each nonnegative integer \(a\) set
\begin{equation}
 \mathsf P_{a,L}(s)=a!\sum_{j=0}^{\lfloor a/2\rfloor}
       \frac{(sL/2)^{a-2j}(L/4)^j}{(a-2j)!j!}.
 \label{eq:moving-kernel-polynomial}
\end{equation}
The differentiated Gaussian formula gives
\(|\partial_z^ak_T|\le\mathsf P_{a,L}(|z-u|)|k_T|\).
In particular the direct first-time factor is
\(L/2+L^2|z-u|^2/4\), by \eqref{eq:time-kernel}; this is a bound for
an actual derivative channel.

The integrated local moments at \(B=2\) are
\begin{subequations}\label{eq:moving-local-moments}
\begin{align}
 M_m[L]&=\sqrt2\e^{R+Y^2L/4+4}
       \sum_{k=0}^m\binom{m}{k}W^{m-k}4^k
                  \frac{\Gamma((k+1)/2)}{\sqrt\pi},\\
 M_m^t[L]&=(L/2)M_m[L]+(L^2/4)M_{m+2}[L].
\end{align}
\end{subequations}
Indeed \(\abs q\le R+\abs v\),
\(\max(\abs q,\abs{z-u})\le W+\abs v\), and
\(\abs v\le v^2/(8T)+2T\). After actual rescaling, the normalized
\(k\)-th Gaussian moment is
\(\sqrt2(8T)^{k/2}\Gamma((k+1)/2)/\sqrt\pi\).
The factor \(T^{-1/2}\) has canceled by integration; now using
\(T\le2\) proves the formulas. With
\(f_L=\exp(Y^2(L-1)/4)\), simultaneously for every \(m\),
\begin{equation}
 M_m[L]=f_LM_m[1],\qquad
 M_m^t[L]\le L^2f_LM_m^t[1].
 \label{eq:moving-local-ratios}
\end{equation}
The finite Bell remainder \(\mathsf B_{p+1}\) in
\cref{lem:weighted} retains degrees through \(2p+2\) for field and
\(2p+4\) for direct time. The complete initial population
\eqref{eq:population} is still
\(C_{\mathrm{pop}}gX^{1/4}\e^{|q|}\),
\(C_{\mathrm{pop}}=4(4\pi)^{-1/4}\). Its integration has no additional
inverse-time cost beyond \eqref{eq:moving-local-moments}.
Radius-\(1/2\) Cauchy disks in the original unit buffers cost at most
\(48\) for spatial orders zero through three.

\subsection{The separate evolved-population ledger}
\label{subsec:moving-population}
Put \(v_n=\log(a_X/n)\), \(a_0=\pi/8\) and
\(c_+=\max((Y-1)/2,0)\). The original slope gates imply
\[
 |\e^{T\lambda_{\epsilon n}^2}|
      \le\e^{2a_0^2+1}\e^{-v_n^2/(4L)}.
\]
The shell proof of \cref{app:population}, including the final weakly
damped cutoff shells, gives \(\Pi_m(1/L,2)\) from \eqref{eq:Pi}, with
\begin{equation}
 W_m(1/L)=2^m+\frac{3^m}{2}\sum_{k=0}^m\binom{m}{k}
                    (8L)^{(k+1)/2}\Gamma((k+1)/2).
 \label{eq:moving-shell-moments}
\end{equation}
Every term then gives
\begin{equation}
 \Pi_m(1/L,2)\le
   \e^{2c_+^2(L-1)}L^{(m+1)/2}\Pi_m(1,2).
 \label{eq:moving-evolved-ratio}
\end{equation}
Keep the unchanged shifted moment majorant
\[
 Z_m(2)=m!\sum_{j=0}^{\lfloor m/2\rfloor}
               \frac{2^j(W+4)^{m-2j}}{j!(m-2j)!}.
\]
The whole shifted \(m\)-th moment sum is bounded by
\(Z_m(2)\Pi_m(1/L,2)gX^{1/4}\). For \(m\le2p+5\), its logarithmic
cost is \(O_{r,b}(L\log(L+2))\), by the same factorial bounds as
\cref{sec:growing}. The degree-dependent power in
\eqref{eq:moving-evolved-ratio} must remain. The real initial population
was already integrated to prove the truncation error; this evolved-mode
bound pays retained-model reconstruction and delivery, and is not an
extra multiplier on that truncation ledger.

\subsection{Moderate, artificial and remote tails}
\label{subsec:moving-tails}
Keep \(\theta=1/2\), \(\delta=1/4\), \(t_R=1+R/4\).
At \(\abs q\ge V\), the same refined split proves
\[
 -(q-\Rea d)^2/8+\abs q
       \le-q^2/16-q^2/32+8t_R^2.
\]
Using \(|z-u|\le2|q|\) and integrating the remaining Gaussian gives
\begin{subequations}\label{eq:moving-tail-moments}
\begin{align}
 \mathsf D_m[L]&=\frac{\sqrt L}{\sqrt{4\pi}}
      \e^{Y^2L/4+8t_R^2}2^m32^{(m+1)/2}\Gamma((m+1)/2),\\
 \mathsf D_m^t[L]&=(L/2)\mathsf D_m[L]+(L^2/4)\mathsf D_{m+2}[L],\\
 \mathsf J_m[L]&=\frac{\sqrt L}{\sqrt{4\pi}}\e^{Y^2L/4}
       \sum_{k=0}^m\binom{m}{k}Y^{m-k}16^{(k+1)/2}
                                      \Gamma((k+1)/2),\\
 \mathsf J_m^t[L]&=(L/2)\mathsf J_m[L]+(L^2/4)\mathsf J_{m+2}[L].
\end{align}
\end{subequations}
The \(\mathsf D\) moments multiply
\(\e^{-V^2/16}=\e^{-8L^2}\); the remote true-data \(\mathsf J\)
moments multiply \(\e^{\pi X/8-X^2/256}\).
Only the remote term uses the complete global source mass. The moderate
term retains \(X^{9/4}\) from the height-sensitive source, and artificial
rays contain only the finite polynomial completion with its full
\(X^{1/4}\) population.

With the actual finite positive polynomials
\(\mathsf B_{p+1},\mathsf H_p\) in \eqref{eq:majorant-polynomials}, set
\begin{subequations}\label{eq:moving-truncation-constants}
\begin{align}
 Q_p[L]&=48\{C_{2p}M_0[L]
                    +C_{\mathrm{pop}}\e\,\mathsf B_{p+1}[M[L]]\}
                      \nonumber\\
 &\quad+C_{2p}M_0^t[L]
                    +C_{\mathrm{pop}}\e\,\mathsf B_{p+1}[M^t[L]],\\
 A_{\mathrm{mod}}[L]&=640\max(\mathsf D_0[L],\mathsf D_0^t[L]),\\
 A_{\mathrm{art},p}[L]&=C_{\mathrm{pop}}
           \max\{\mathsf H_p[\mathsf D[L]],\mathsf H_p[\mathsf D^t[L]]\},\\
 A_{\mathrm{far}}[L]&=16\max(\mathsf J_0[L],\mathsf J_0^t[L]).
\end{align}
\end{subequations}
The real initial mismatch has the unchanged all-positive-integer-order
constant
\[
 C_{2p}=\frac{2^{7/4}}7\Gamma(p+1/2)1.1^{2p+1}(8\pi)^{p+3/4},
\]
and the separate retained-size gate \(K\le a_{\min}^2\) remains.
Since every \(\mathsf D_m[L]=\sqrt L f_L\mathsf D_m[1]\), and the
same holds for \(\mathsf J\), a common multiplier for these four
truncation prefactors is
\begin{equation}
 \rho_{\mathrm{tr}}(L)=L^{5/2}\exp(Y^2(L-1)/4).
 \label{eq:moving-truncation-ratio}
\end{equation}
The time versions cost at most another \(L^2\). Positivity of the finite
polynomial sums proves this multiplier uniformly in every retained
order. It does not bound all evolved moments in
\eqref{eq:moving-evolved-ratio}, or all complex provider magnitudes.

Insert the exact initial identity into \eqref{eq:heat-convolution},
cancel the bounded \(Q\)-integral and subtract the whole-line polynomial
transforms exactly as in \cref{thm:fixed}. The result before absorption is
\begin{align}
 \lVert h-F_X^{[p,2p]}\rVert_{\gamma;X;\mathrm{mov}}^{(3,1)}
 \le C_\chi\bigl[&Q_p[L]X^{-(p+3/4)}\nonumber\\
 &+49\{A_{\mathrm{mod}}[L]X^{9/4}
                 +A_{\mathrm{art},p}[L]X^{1/4}\}\e^{-8L^2}\nonumber\\
 &+49A_{\mathrm{far}}[L]X^{-7/4}
                                  \e^{\pi X/8-X^2/256}\bigr].
 \label{eq:moving-unabsorbed}
\end{align}
This pays source mismatch, local weighted remainder, moderate true data,
artificial polynomial tails and remote true data separately. No real-part
operation is moved across a complex-kernel integral; no true gamma or
Riemann--Siegel function is extended along artificial polynomial rays.

For each frozen chart, \(1/L>0\). Dominated differentiation of the fixed
bounded integral and the finite Gaussian transforms proves
\((\partial_T-\partial_z^2)F=0\) and its first spatial derivative on a
positive-time neighborhood. This exact identity does not discard any of
the preceding comparison errors.

\subsection{Continuous onset and half-target slack}
\label{subsec:moving-onset-proof}
The envelope proof in \cref{sec:calibration} uses time only through
\(0<T\le2\): its logarithmic discrepancy is multiplied by \(T/16\).
Thus \eqref{eq:explicit-envelope-lower} and
\(c_{\mathrm o}=1\) extend to this entire range, giving
\[
 \calE_{\mathrm o}(-T,z)\ge|G(z)|
                 \e^{-TL^2/16+\pi^2T/64}.
\]
Natural-to-outer conversion therefore still costs at most
\(\e^{L^2/8}\). The disk, slope and complete-source calibration constants
are independent of the lower time endpoint. Every structural gate in
\cref{thm:linear-onset} follows from \(L\ge64\nu\); none is replaced by
a fixed-order tail absorption gate.

For \(L\ge1024\nu^2\), \(L\ge16384\), \(Y<\nu\) and the continuous
logarithm estimate used in \cref{subsec:linear-onset} gives
\(\log L\le L/32\). Hence
\begin{equation}
 \log\rho_{\mathrm{tr}}(L)
 \le\nu^2L/4+(5/2)\log L
 \le L^2/4096+5L/64\le L^2/2048.
 \label{eq:moving-rho-slack}
\end{equation}
The last step needs only \(L\ge320\).
Add this loss to the actual additive envelopes in the proof of
\cref{thm:linear-onset}, retaining the half-target factors \(8,392\):
\begin{align}
 \Lp(8C_{\mathrm o}Q_p[L]),\quad
 \Lp(392C_{\mathrm o}A_{\mathrm{art},p}[L])
                      &<193L^2/2048,\nonumber\\
 \Lp(392C_{\mathrm o}A_{\mathrm{mod}}[L]),\quad
 \Lp(392C_{\mathrm o}A_{\mathrm{far}}[L])
                      &\le65L^2/2048.
 \label{eq:moving-prefactor-slack}
\end{align}
Here \(C_{\mathrm o}=C_\chi/c_{\mathrm o}=3^{\nu-1}\).
The first half-target gap satisfies
\[
 (p+3/4)\log X-3L^2/8\ge7L^2/64=224L^2/2048,
\]
leaving strict slack \(31L^2/2048\). The moderate and artificial
right sides of \eqref{eq:growing-gates} cost at most
\(577L^2/2048\) and \(321L^2/2048\), respectively, each less than
\(61L^2/8\). For the remote gate the same earlier proof gives
\(\pi X/8<X^2/512\) and
\(X^2/512>9L^2/8>(3/8+65/2048)L^2\).
Thus all four half-target inequalities close. After outer conversion,
\eqref{eq:moving-unabsorbed} uses at most \(\tfrac12\e^{-L^2/4}\).

Only the floor bounds \(L/2-1\le p\le L/2\) enter these continuous
majorants. The complete population bound and exact adjacent-cutoff
transfer hold through every later \(p,N\) jump. No monotonicity of the
actual constants or continuity of the selected model in \(X\) is assumed.
The dyadic sufficient condition follows from
\((2L_{\mathrm{mov}}+8)\log2-\log(4\pi)>L_{\mathrm{mov}}+1\).
The moving sufficient log-height is \(16\nu\) times the old \(64\nu\)
bound; this is the cost of this proof, rather than a necessary threshold.

\subsection{Recomposition of all delivered-error channels}
\label{subsec:moving-allocation}
Retain the original five evaluation categories plus scale. The center
budget and its comparable finite rational choice are
\begin{equation}
 \varepsilon_X=\frac{1}{2C_\chi}\e^{-3L^2/8},\qquad
 \frac{1}{8C_\chi}\e^{-3L^2/8}<e<
                  \frac{1}{4C_\chi}\e^{-3L^2/8}.
 \label{eq:moving-budget}
\end{equation}
Use factor-controlled rational upper enclosures in the finite allocation
of \cref{subsec:finite-allocation}, making the following explicit
substitutions.

\begin{enumerate}[leftmargin=*,itemsep=.55em]
\item \textbf{Arithmetic data and scale.}
Use \eqref{eq:moving-local-moments} in place of \(M,M^t\), so that
\(\mathfrak M_m[L]=48M_m[L]+M_m^t[L]\).
Recompute \(\mathcal D_p,\mathcal D_p^{(1)},\mathcal C_p,
\mathcal S_p\) in \cref{sec:evaluation}. The exact coefficient
sensitivity \eqref{eq:coefficient-sensitivity} and full population
in \eqref{eq:arithmetic-error} remain. Complete-square carrier bounds
become \(\exp(R+2+Y^2L/4)\), and
\eqref{eq:moving-evolved-ratio} pays the distinct evolved reconstruction.
Keep the slope neighborhood, all \(O(X)\) phase multipliers, unrestricted
scale exponents and all \(N\) summands. The recomputed bounds require
\(E_{\mathrm{ar}}\le\varepsilon_X/20\) and
\(\varepsilon_gA_0\le\varepsilon_X/2\).

\item \textbf{Provider-domain kernels and retained data.}
Keep the original \(B_Q,B_Q^{\mathrm{loc}},a_*\), real branches and
holomorphic extensions of \cref{sec:quadrature,sec:evaluation}; replace only
their kernel bounds by
\begin{subequations}\label{eq:moving-provider-kernels}
\begin{align}
 K_E[L]&=\frac{\sqrt L}{\sqrt{4\pi}}
       \e^{L(Y+3V/4)^2/4}
       \left[\max_{a\le3}\mathsf P_{a,L}(W+2V)
                           +\mathsf P_{2,L}(W+2V)\right],\\
 K_{\mathrm{loc}}[L]&=\frac{\sqrt L}{\sqrt{4\pi}}
       \e^{L(Y+1)^2/4}
       \left[\max_{a\le3}\mathsf P_{a,L}(W+V+4)
                           +\mathsf P_{2,L}(W+V+4)\right].
\end{align}
\end{subequations}
Set \(B_E[L]=B_QK_E[L]\),
\(B_{\mathrm{loc}}[L]=B_Q^{\mathrm{loc}}K_{\mathrm{loc}}[L]\), and
\[
 A_0=C_{\mathrm{pop}}X^{1/4}\mathsf H_p[\mathfrak M[L]]
                            +2VB_{\mathrm{loc}}[L].
\]
The normalized weighted-table and pole-safe contour constructions retain
every \(C_k\). Either \cref{prop:sharper-rs,prop:smaller-contour} or the
old provider applies. Its coefficient gain is not applied to an entire
transfer-containing bound. All contour, data, node, weight and \(\pi\)
errors, including the constant \(+1\) in \eqref{eq:Ck-delivery}, remain.
Recompute \cref{eq:node-data-error,eq:RS-data-budget} with the new local
kernel, choosing the old finite tolerances to give
\(E_{\mathrm{RS}}\le\varepsilon_X/20\).

\item \textbf{Exact outer quadrature.}
The same positivity and Chebyshev proof gives
\[
 E_{\mathrm{quad}}\le16VB_E[L]4^{-n}.
\]
Choose \(n\ge1\) by the strict rational test
\(16\overline V\overline B_E[L]4^{-n}<e/80\).
The test terminates and pays \(E_{\mathrm{quad}}\le\varepsilon_X/20\).
For the unchanged ellipse of parameter \(2\),
\begin{equation}
 \Lp B_E[L]=O_{r,b}(L^3+L\log(L+2)),\qquad
 n=O_{r,b}(L^3).
 \label{eq:moving-outer-order}
\end{equation}
Its imaginary height is \(O(L)\); the squared kernel displacement is
divided by \(1/L\). This explains the cubic cost. It is a sufficient
rule for this provider, not a lower bound on all quadrature rules.

\item \textbf{Frozen geometry.}
Retain exact branch labels, interior rational guards and
\eqref{eq:geometry-error}, now with the recomputed bound:
\[
 E_{\mathrm{geo}}\le B_{\mathrm{loc}}[L]
  \{\varepsilon_w+2(2V+\varepsilon_w)\varepsilon_u+6\varepsilon_b\}.
\]
The same finite tolerances give \(E_{\mathrm{geo}}\le\varepsilon_X/20\).
No new real edge-equality decision is introduced. Reference intervals
remain inside true branches, and mapped node/weight effects are charged
once.

\item \textbf{Remaining circuit and rounded outputs.}
Use \cref{eq:roundoff-updates,eq:channel-formulas} with the actual new
kernel and carrier bounds, safe reciprocals and full accumulation count.
Earlier preprocessing errors are paid once. Refine primitive and
operation tolerances until \(E_{\mathrm{round}}\le\varepsilon_X/20\).
For the literal endpoint tag, its time-input split in
\cref{subsec:moving-input} is included in this category. Exact rational
queries have no time-input error.
\end{enumerate}
Consequently
\(\varepsilon_gA_0+(1+\varepsilon_g)\sum E_i\le\varepsilon_X\).
Together with the half-target comparison, this proves
\eqref{eq:moving-delivery}, subject to the input construction below.
The exact frozen residual does not remove any category. Independently
rounded channels have only the sum-of-channel heat-defect allowance;
its spatial derivative still requires an additional mixed channel.

\subsection{Strict rational queries and the literal endpoint tag}
\label{subsec:moving-input}
Keep the exact positive dyadic \(X,r,b\), supplied \(p,N\), and strict
rational descriptor enclosures from \cref{subsec:input-class}:
\begin{equation}
 2p<L^-\le L\le L^+<2p+2,\qquad
 N<a^-\le\sqrt{X/(4\pi)}\le a^+<N+1.
 \label{eq:moving-descriptors}
\end{equation}
Include strict gate/half-target witnesses and exact rational displacement
\(d=z-X\) inside the original buffered chart. An ordinary query supplies
exact rational \(T\le2\) and the additional finite witness
\begin{equation}
 TL^->1.
 \label{eq:moving-strict-time-witness}
\end{equation}
Invalid finite witnesses are rejected. Every rational \(T>1/L\) has a
witness after refinement; the search has nonuniform separation
\(TL-1\). A rational query below the lower endpoint is rejected.
No equality oracle for \(T=1/L\), descriptor boundaries or arbitrary
computable reals is supplied. The mathematical comparison is uniform
on the closed interval, whereas this ordinary interface has a strict
input class.

The endpoint has a separate finite interface: the literal tag
\(\texttt{T=1/log(X/(4pi))}\) defines one computable number from the
already exact \(X\). It is not an equality test on an arbitrary input.
\begin{enumerate}[leftmargin=*,itemsep=.5em]
\item Refine positive bounds until \(L^+/L^-<2\) and \(L^->1\).
On the coarse interval \(I=[1/L^+,1/L^-]=[T_-,T_+]\),
\(T_->1/(2L)\) and \(T_+<2/L\le2\).

\item Freeze all delivered coefficients, nodes, weights and data and
regard their center-normalized finite polynomial/Gauss model as a function
of real \(T\).
It solves heat exactly, so
\(\partial_T\partial_z^a=\partial_z^{a+2}\) and
\(\partial_T^2=\partial_z^4\). Bound the time derivative of every
requested channel on \(I\) by one rational \(B_T\): use the finite spatial
formula \eqref{eq:channel-formulas} through order five for the arithmetic
part, and \eqref{eq:moving-kernel-polynomial} through order five with
inverse time \(1/T_-\) for every kernel node. Twice the maximum of these
center-normalized spatial orders zero through five pays the combined
center channel norm.
Positive coefficient/node bounds and unrestricted exponent ranges make
this bound finite and computable. This pays \(F_{TT}\), which is needed
when the first-time output is approximated.

\item Refine the endpoint enclosure to width
\(\delta_T<e/[640(1+B_T)]\), evaluate at its rational midpoint, and charge
\(B_T\delta_T<e/640\) as time-input error to \(E_{\mathrm{round}}\).
Use safe inverse time \(1/T_-\) throughout that circuit. Allocate the
remaining new circuit and quantization error below \(e/160\); together
they are below \(e/80<\varepsilon_X/20\). Earlier coefficient, RS-data
and geometry payments are not charged again.
\end{enumerate}
The midpoint may be slightly below \(1/L\). Only the fixed delivered
model is evaluated there; truncation is invoked at the exact tagged time.
The coarse interval is bounded away from zero, so refinement terminates.
For fixed charts the derivative bound has log-magnitude
\(O_{r,b}(L\log(L+2))\), by the finite moment/path bounds with
\(T_->1/(2L)\). Thus the tag fits \eqref{eq:moving-precision}, counted
in the actual input and tolerance lengths. It adds no general uncertain
coordinate theorem or mixed-residual certification. This completes the
input and delivery proof of \cref{thm:moving-endpoint-delivery}.

\subsection{Delivered precision, internal work and remaining scope}
\label{subsec:moving-cost}
On the real/local strip,
\(\Lp B_{\mathrm{loc}}[L]=O_{r,b}(L\log(L+2))\), as for arithmetic
conditioning and positive circuit paths. The allocations above therefore
require tolerance at worst
\(\exp[-3L^2/8-O_{r,b}(L\log(L+2))]\), proving
\eqref{eq:moving-precision}. Full summation and \(O(L^3)\) outer nodes
add their logarithmic counts. The off-real ellipse magnitude pays
quadrature truncation; it must not be substituted for the real-node
precision bound. The RS contour order remains
\(O_{r,b}(L^2+L\log(L+2))\); its direct scalar composition count becomes
\(O(L^6)\), in addition to the full \(O(Np)\) arithmetic work.

Provider-internal root/weight separation depths can greatly exceed
requested delivered precision. For the deliberately loose bound parallel
to \cref{thm:bit-cost}, let \(\mathcal B\) count actual finite input and
witness lengths, retain the descriptor separation \(s_\sigma\), and put
\begin{equation}
 \begin{gathered}
 s_T=\left\lceil\log_2^+\frac1{TL-1}\right\rceil
       \quad\hbox{for searched ordinary rational witnesses},\\
 s_T=0\quad\hbox{for the literal endpoint tag},\qquad
 Q_{\mathrm{mov}}=2+\mathcal B+P+s_\sigma+s_T+\lceil L\rceil^3.
 \end{gathered}
 \label{eq:moving-work-parameter}
\end{equation}
Verification of supplied witnesses is charged through their actual
length in \(\mathcal B\); \(s_T\) records their nonuniform acquisition
cost when they are found by refinement. Then the same degree accounting
gives effective chart constants with
\begin{equation}
 T_{\mathrm{bit}}\le C(r,b)(N+1)Q_{\mathrm{mov}}^{40},\qquad
 M_{\mathrm{bits}}\le C(r,b)Q_{\mathrm{mov}}^{20}.
 \label{eq:moving-work}
\end{equation}
Indeed \(n,q,p,K\) and the gamma shift order are \(O(Q_{\mathrm{mov}})\),
inverse time is \(O(Q_{\mathrm{mov}})\), and real-node magnitude bit
lengths remain polynomial. Gauss isolation depth is
\(O(Q_{\mathrm{mov}}^4)\), mapped-weight lengths
\(O(Q_{\mathrm{mov}}^5)\), gamma work
\(O(Q_{\mathrm{mov}}^{24})\) per call and
\(O(Q_{\mathrm{mov}}^{25})\) over all such calls. Remaining
\(O((N+1)Q_{\mathrm{mov}}^{35})\) work fits the displayed bound, with
the same streaming argument. The original \(\lceil L\rceil^2\) work
parameter is not silently reused.

These bounds do not claim measured speed, sharp complexity or complete
software validation. The theorem excludes \(T=0\), a fixed-height limit
as \(T\downarrow0\), uniform unbounded radii, physical-observable
transfer and any conclusion about zeros. Strong negative precision
through zero remains open. The positive signed-crossover atlas in
\cite{ansatzcomp} has its own weaker field allowances and does not
extend the present strong precision target. The old \([1,2]\) theorem,
its \(64\nu\) onset and original providers remain valid fallbacks.
